\chapter[Train Track Case Analysis: Track 12]{Train Track Case Analysis: Track 12}
\label{app:case_analysis_track12}

This is Track 17:

\section{Case A: Triangles Fixed}\label{CaseA Track 12}
Let us assume the triangles are fixed. Hence, both $\fbd$ and $\fbdb$ both pass over $d$/$\dbar$ after an even number of letters. We shall call this \textit{Assumption A}. Then $\fbd$ can begin with $b$ or $r$, while $\fbdb$ can begin with $s$, $i$, $p$, or $o$.

\subsection{BKG}\label{BKG}
Suppose $\fbd$ began with $b$, and $\fbdb$ began with $s$.

This is BKG1.

Here in order to meet Assumption A, $\fbdb$ must eventually traverse $p^+$ before crossing $\dbar$. Before that it would possibly have twisting with $o$ or $i$.

Here we claim that the dotted lines must join in the middle. Indeed, if they did not, then in order to not put either $\fbd$ or $\fbdb$ in Situation $\beta$, this can only occur with $\fbdb$ going below $\fbr$, or $\fbp$ going above $\fbd$. In the former case, this causes $\fbp$ to trace, as shown in BKG2. In the latter case, the dotted lines will be in Situation $\alpha$, as shown in BKG3.

Thus the dotted lines must join, so we have BKG4.

Now let us assume that there were no initial twists by $\fbdb$. Then $\fbr$ will trace, as shown in BKG5.

So let us assume instead that $\fbdb$ has initial twists with $o$. 

But this means $\fbo$ will trace, as shown in BKG6.

% Then it must be
% \[
% \fbdb = s^- o^+ s^+...
% \]

% This is BKG6. Notice here that $\fbr$ forces either
% \[
% \fbo = s^- o^\times
% \]
% or
% \[
% \fbo = s^- \dbar b^- r^- d p^- s^*...
% \]

% The latter case is shown in BKG7. Here the dotted lines are forced to join, otherwise $\fbd$ will violate Rule F1.

% Now consider $\fbb$ and $\fbp$. $\fbb$ can only begin with $r^\circ$ or $r^-$. Suppose it were the former case. Then we must have
% \[
% \fbp = \dbar r^+ b^\circ
% \]

% as continuing with either $b^+$ or $b^-$ would cause it to cross $d$ and immediately trace. Now this forces $\fbi$ to begin with $p^-$, otherwise $\fbs$ would begin with $p^+$ and trace:
% \[
% \fbs = p^+ \dbar r^+ b^+ r^- d s^\times
% \]

% Thus we must have
% \[
% \fbi = p^- s^- o^*...
% \]

% This immediately forces $\fbo$ to end at $i^\circ$. Otherwise, it either circles in the outer perimeter and trace, or it re-enters and then due to the position of $\fbi$, will trace.

% Now both $\fbs$ and $\fbi$ must end. Otherwise, $\fbs$ can only begin with $p^-$ which would cause it to immediately trace. This gives us Candidate 1, shown in BKG8.

% Now suppose
% \[
% \fbb = r^- d p^*...
% \]

% as shown in BLG9. Here $\fbr$ must end otherwise it would trace after crossing $\dbar$, regardless it it continued with $s^+$ or $s^-$.

% Also here $\fbp$ must end, or it would immediately trace after crossing.

% Here $\fbb$ can end at $p^\circ $ or continue with $p^-$ (to the right of $\fbs$). Suppose it were the former. Then we have BLG10. Here $\fbi$ must continue with $b^-$ and to the right of $\fbs$. Otherwise, $\fbs$ would eventually trace. So we have
% \[
% \fbi = ...b^- r^- d p^- s^- o^*...
% \]

% which means immediately that $\fbs$ must end or it would eventually trace; and $\fbo$ must also immediately end or it would either circle the outer perimeter and trace, or re-enter and follow $\fbi$ and thereby trace. This gives Candidate 2, shown in BKG11.

% Now suppose $\fbb$ continued with $p^-$:
% \[
% \fbb = r^- d p^- s^- o^*...
% \]

% This means that $\fbi$ must begin with $p^-$. Otherwise, $\fbs$ would trace:
% \[
% \fbs = p^+ \dbar r^+ d s^\times
% \]

% So we have
% \[
% \fbi = p^- s^- o^*...
% \]
% and
% \[
% \fbs = p^\circ
% \]
% otherwise it would immediately trace.

% This is BKG12. 

% Here $\fbo$ again must end due to the position of $\fbi$, just like before.

% Here $\fbi$ must continue with $o^+$ and to the left of $\fbb$. Otherwise, $\fbb$ would be absorbed:
% \[
% \fbb = ...o^- s^+ p^+ \Abso{\fbi}{\fbs}
% \]

% So we must have
% \[
% \fbi = ...o^+ s^+ p^+ \dbar r^+ b^*...
% \]

% and here $\fbb$ must end immediately or it would trace due to the position of $\fbi$, and in turn $\fbi$ must end immediately. This is Candidate 3, shown in BKG13.

% Going back to the situation depicted in BKG4. Suppose now that $\fbdb$ had initial twists with $p$. Then it must be:
% \[
% \fbdb = s^+ p^-...
% \]
% and here if the next letter were $i$, then $\fbi$ would trace. But if it were $s$, then $\fbs$ would trace.

This is a contradiction.

\subsection{CNM}
Suppose $\fbd$ began with $b$, and $\fbdb$ began with $i$.

The same argument used in BKG \ref{BKG} to show that the dotted lines in the middle must join applies here.

This is CNM1.

Here if $\fbdb$ were to begin with $i^+$ then it must continue with $p^+$ otherwise $\fbi$ would immediately trace. Moreover, $\fbdb$ must join with $\fbd$ here or $\fbd$ would be in Situation $\beta$. So in this scenario we have CNM2.

Here $\fbs$ must begin with $p^+$ otherwise $\fbi$ will immediately trace. So it continues:
\[
\fbs = p^+ \dbar r^+...
\]
where it must be $r^+$ otherwise it would cause $\fbp$ to trace after crossing.

Here $\fbs$ cannot go into 2, otherwise it will escape to the outer perimeter and eventually trace due to the position of $\fbr$. So it goes into 1, which is CNM3.

Here $\fbb$ can either begin with $r^\circ$ or $r^-$. Suppose it were the former. Then we have $\fbp$ and $\fbs$ both continuing with $r^+$, and $\fbs$ must then continue with $b^-$. Otherwise, $\fbp$ would continue with $b^+$ and then trace:
\[
\fbp = ...b^+ r^- d p^\times
\]

So we have
\[
\fbs = ...b^- r^- d...
\]
This is CNM4. But here $\fbs$ escapes to the outer perimeter again and will eventually trace.

So suppose instead that $\fbb$ began with $r^-$, as in CNM5.

Here it can go into 1 or 2. Suppose it went into 1. Then
\[
\fbb = r^- d s^- o^*...
\]
where it must traverse $s^-$ as to not cause $\fbr$ to traverse $s^+$ and thereby trace.

This is CNM6.

Here again we must have
\[
\fbs = ...b^- r^- d...
\]
and $\fbs$ will escape to the outer perimeter and trace, as shown in CNM7.

So suppose $\fbb$ went into 2. Then we have the situation depicted in CNM8. Here $\fbi$, $\fbs$, and $\fbb$ are stuck in a three-way counterclockwise path mill between $\fbb$, and $\fbp$. This is a contradiction.

So going back to the situation depicted in CNM1, suppose $\fbdb$ began with $i^-$, then
\[
\fbdb = i^- o^*...
\]
as it cannot cross or it would violate Assumption A. But now $\fbo$ will follow it and trace.

This is a contradiction.

\subsection{LGU}
Suppose $\fbd$ began with $b$, and $\fbdb$ began with $p$. Then $\fbdb$ must begin with $p^-$, as shown in LGU1; and then it must traverse either $i^- \dbar$, or $s^- \dbar$. In either case, this would cause the dotted lines to be in Situation $\alpha$, as shown in LGU2, and LGU3.

This is a contradiction.

\subsection{CXH}
Suppose $\fbd$ began with $b$, and $\fbdb$ began with $o$. But this forces $\fbo$ to also begin with $o$ which is impossible.

This concludes the case when we assume $\fbd$ begins with $b$. Now we assume $\fbd$ begins with $r$, and $\fbdb$ can begin with $p$, $i$, $s$, or $o$, as before. Again Assumption A must be upheld.

\subsection{SXM}
Suppose $\fbd$ began with $r$, and $\fbdb$ began with $p$. Then we have the situation depicted in SXM1.

Here there are four possibilities for $\fbdb$ to traverse before crossing without violating Assumption A. We have
\[
\fbdb = p^- i^- \dbar...
\]
\[
\fbdb = p^- s^- \dbar...
\]
\[
\fbdb = p^- s^- o^- (s^- o^+)^k s^+...
\]
\[
\fbdb = p^- i^- o^- (i^- o^+)^k i^+...
\]
where $k = 0,1,2,...$ is a number of possible twisting.

Let us first assume the first case. Notice that in the middle, $\fbb$ must go above the pair $\fbi$, and $\fbs$. Otherwise, the pair would traverse $\dbar r^+ b^+ d$, and then either $\fbs$ or $\fbi$ would trace. So we have
\[
\fbb = d p^- i^- o^*...
\]
as shown in SXM2.

Here if $\fbo$ were to go above $\fbd$, then it would force $\fbdb$ into Situation $\beta$, But if it went below $\fbd$, then it would eventually trace.

So let us assume the second case, that 
\[
\fbdb = p^- s^- \dbar...
\]

This is SXM3, where again $\fbb$ must go above $\fbdb$ for the same reason as before, and the dotted lines must join in the middle otherwise one of them will immediately enter Situation $\beta$.

Here $\fbo$ must continue with $i^+$ and to the left of $\fbb$ otherwise $\fbb$ would trace after crossing $\dbar$. But now $\fbo$ will be absorbed:
\[
\fbo = p^- s^- \dbar b^- r^- d p^- \hat{i}^+ p^+ \dbar r^+ b^+ d s^+ p^+ \Abso{\fbo}{\fbdb}
\]

So suppose we have the third case:
\[
\fbdb = p^- s^- o^- (s^- o^+)^k s^+...
\]

(in the following analysis for simplicity we draw these cases with $k=1$. For more twisting the arguments will remain the same with possibly more redundant twisting inserted into each image path)

This forces
\[
\fbp = s^- o^*...
\]

Again the dotted lines must join and $\fbb$ traverses

\[
\fbb = d p^- i^*..
\]
and also
\[
\fbo = p^- i^*...
\]

And again $\fbo$ must continue with $i^+$:
\[
\fbo = p^- i^+ p^+ \dbar r^+ b^+ d s^*...
\]

This is SXM5.

Here $\fbb$ must end to give
\[
\fbb = d p^- i^\circ
\]
as either continuing with $i^+$ or $i^-$ would cause it to immediately trace after crossing $\dbar$.

This is SXM6.

Here consider $\fbr$. It can either begin with $b^\circ$ or $b^+$. Suppose it were the former. This is SXM7. 

Here if $\fbi$ continued with $r^\circ$ or $r^+$, forcing $\fbs$ to continue with $r^+$, then $\fbs$ would trace immediately after crossing $d$. On the other hand, if $\fbi$ continued with $r^-$, then it would trace:
\[
\fbi = ...r^- b^+ d p^- i^\times
\]

This is a contradiction.

So suppose $\fbr$ began with $b^+$. Then it must do so going to the left of $\fbi$, as otherwise it would either be absorbed, or $\fbi$ and $\fbs$ would again be put in the same situation as before. So we have
\[
\fbr = b^+ d s^- o^*...
\]

This is SXM8.

Here if $\fbi$ continued with $b^\circ$ or $b^+$, forcing $\fbs$ to continue with $b^+$, then $\fbs$ is immediately absorbed. But if $\fbi$ continued with $b^-$ then it will trace:
\[
\fbi = \dbar \hat{b}^- d p^- i^+ p^+ \dbar r^+ b^+ d s^\times
\]

This is a contradiction.

So now suppose it were the fourth case:
\[
\fbdb = p^- i^- o^- (i^- o^+)^k i^+...
\]


But now $\fbo$ will trace as shown in SXM9.

\subsection{GIB}\label{GIB}
Suppose $\fbd$ began with $r$, and $\fbdb$ began with $i$. Then we have the situation depicted in GIB1.

Here, as show, in order to meet Assumption A, $\fbdb$ must eventually traverse $p^+$ right before crossing, possibly first with some twisting.

Here if the dotted lines do not join, then in order to not put either $\fbd$ or $\fbdb$ in Situation $\beta$, the only way this could be done is if $\fbp$ went above $\fbd$; or $\fbdb$ went below $\fbd$.

In the former case, this causes $\fbp$ to trace, as shown in GIB2. In the latter case, the dotted lines will enter Situation $\alpha$, as shown in GIB3.

So the dotted lines must join. But in this case $\fbp$ still trace, as shown in GIB4.

This is a contradiction.

\subsection{TEX}
Suppose $\fbd$ began with $r$, and $\fbdb$ began with $s$. Then we have the situation depicted in GIB1. As in the situation in GIB \ref{GIB}, $\fbdb$ must eventually traverse $p^+$ right before crossing.

Here we can repeat the entire argument in GIB \ref{GIB} to reach a contradiction.

\subsection{LYR}
Suppose $\fbd$ began with $r$, and $\fbdb$ began with $o$. But this is impossible as this forces $\fbo$ to also begin with $o$.

This concludes the case analysis for Case A: Triangles Fixed.

\section{Case B,C,D: Triangles Swapped}\label{CaseBCD Track 12}

Now suppose the triangles are swapped. Then at least one of $\fbd$ and $\fbdb$ must go over an odd number of letters before crossing. We first assume this to be the case for $\fbdb$, that is, $\fbdb$ crosses an odd number of letters before crossing $d$.

Then $\fbdb$ can begin with $r$, $b$; and $\fbd$ can begin with $p$, $i$, $i$, $s$ or $\dbar$.

\subsection{TGU}\label{TGU}
Suppose $\fbdb$ began with $r$, and $\fbdb$ began with $p$. We remind the reader that we are assuming $\fbdb$ must go over an odd number of letters before crossing. This assumption necessitates that eventually $\fbdb$ must traverse $o^-$ right before crossing $d$.

This is TGU1.

Here if $\fbr$ went above $\fbp$ then $\fbr$ will immediately trace; if instead $\fbr$ went below $\fbp$ then $\fbp$ will immediately trace.

This is a contradiction.

\subsection{XNA}\label{XNA}
Suppose $\fbdb$ began with $r$, and $\fbdb$ began with $i$. We remind the reader that we are assuming $\fbdb$ must go over an odd number of letters before crossing. This assumption necessitates that eventually $\fbdb$ must traverse $o^-$ right before crossing $d$.

This is XNA1.

Here $\fbb$ must go above $\fbp$ to not immediately trace; and in between $\fbp$ and $\fbdb$ in order to not trace as well. This forces $\fbdb$ to not have any initial twists or $\fbb$ would be brought along and trace. 

This is XNA2.

Here $\fbdb$ must continue with $p^-$ after crossing to not violate Rule F1.

Here if $\fbs$ began with $b^\circ$ or $b^-$, forcing $\fbi$ to begin with $b^-$ then $\fbi$ will trace:
\[
\fbi = b^- r^- d p^- i^\times
\]

So $\fbs$ must begin with $b^+$:

\[
\fbs = b^+ d i^*...
\]
forcing 
\[
\fbp = d i^*...
\]

This is XNA3.

Here $\fbd$ must begin with $i^-$ otherwise $\fbp$ would trace. Then if $\fbd$ immediately crosses, it would traverse
\[
\fbd = i^- \dbar b^- r^- d p^-...
\]
and then the dotted lines will be in Situation $\alpha$. So we must have
\[
\fbd = i^- o^*...
\]
instead.

But now as XNA4 shows, $\fbo$ will trace.

This is a contradiction.

\subsection{WYS}\label{WYS}
Suppose $\fbdb$ began with $r$, and $\fbdb$ began with $s$. We remind the reader that we are assuming $\fbdb$ must go over an odd number of letters before crossing. 

This is WYS1.

The same arguments used in XNA \ref{XNA} lead to the conclusion that $\fbs$ must begin with $b^+$, but in this instance this means it will immediately trace after crossing, as shown in WYS2.

This is a contradiction.

\subsection{HTI}\label{HTI}
Suppose $\fbdb$ began with $r$, and $\fbdb$ began with $o$. We remind the reader that we are assuming $\fbdb$ must go over an odd number of letters before crossing. This assumption necessitates that eventually $\fbdb$ must traverse $o^-$ right before crossing $d$.

In this case, $\fbp$ is forced to immediately trace after crossing, as shown in HTI1.

This is a contradiction.

\subsection{SEZ}\label{SEZ}
Suppose $\fbdb$ began with $r$, and $\fbdb$ began with $\dbar$. We remind the reader that we are assuming $\fbdb$ must go over an odd number of letters before crossing. This assumption necessitates that eventually $\fbdb$ must traverse $o^-$ right before crossing $d$.

Here $\fbp$ must go under $\fbd$ or it would immediately trace.

We have SEZ1.

Here $\fbp$ can continue with $i$ or $o$. Suppose it were the former. Then this forces 
\[
\fbr = i^- o^*...
\]

to prvent $\fbp$ from tracing. Also, we must have
\[
\fbi = b^- r^- d...
\]
to prevent $\fbs$ from tracing after crossing $d$.

We have the situation depicted in SEZ2.

The position of $\fbi$ here forces $\fbdb$ to not have any initial twists, so we have
\[
\fbfb = r^- d...
\]
and
\[
\fbo = r^- d...
\]
and here $\fbo$ must go above $\fbd$, otherwise if $\fbd$ went in between $\fbo$ and $\fbi$ then it would be in Situation $\beta$; while if it went above $\fbi$, it would cause $\fbi$ to trace. So now
\[
\fbo = r^- d p^*...
\]
and here it must continue with $p^-$ or else $\fbb$ would trace.

But now we have the situation depicted in SEZ3, where it is clear that $\fbo$ will trace.

So suppose instead $\fbp$ continued with $s$:
\[
\fbp = d s^*...
\]
and if $\fbr$ also began with $s$ here we would have SEZ4. Here the same argument as before shows the dotted liens must join. Here $\fbi$ and $\fbo$ must continue with $p^-$ otherwise $\fbb$ would trace after crossing; but then one of $\fbo$ or $\fbi$ will trace.

If $\fbr$ instead began with $i$, then this is the situation depicted in SEZ5. Here again $\fbi$ and $\fbo$ must continue with $p^-$ otherwise $\fbb$ would trace after crossing; but then one of $\fbo$ or $\fbi$ will trace.

This is contradiction.

This concludes the situation where we assume $\fbdb$ goes over an odd number of letter before crossing. Now we assume $\fbdb$ goes over an even number of letters before crossing, which means $\fbd$ must go over an odd number of letters before crossing.

Hence, $\fbdb$ can begin with $b$ or $r$, whicle $\fbd$ can begin with $p$, $i$, $s$, $o$.

\subsection{PBH}
Back to the end of Case B,C,D \ref{CaseBCD Track 12}.
Suppose $\fbdb$ began with $b$, and $\fbdb$ began with $p$. We assume $\fbd$ must go over an even number of letters before crossing, which necessitates that $\fbdb$ must go over an odd number of letters before crossing.

Here we must have
\[
\fbdb = b^- r^- d...
\]
and thus
\[
\fbo = b^- r^- d...
\]
as shown in PBH1.

Here no matter where $\fbr$ continues it will trace.

This is a contradiction.

\subsection{CVF}\label{CVF}
Back to the end of Case B,C,D \ref{CaseBCD Track 12}.
Suppose $\fbdb$ began with $b$, and $\fbdb$ began with $i$. We assume $\fbd$ must go over an even number of letters before crossing, which necessitates that $\fbdb$ must go over an odd number of letters before crossing.


Here the only way $\fbb$ can not trace is to go in between $\fbdb$ and $\fbi$. But then $\fbi$ and $\fbs$ cross, and one of them will immediately trace on the other side.

This is a contradiction.


\subsection{BRW}
Back to the end of Case B,C,D \ref{CaseBCD Track 12}.
Suppose $\fbdb$ began with $b$, and $\fbdb$ began with $s$. We assume $\fbd$ must go over an even number of letters before crossing, which necessitates that $\fbdb$ must go over an odd number of letters before crossing.

The same argument of CVF \ref{CVF} applies to give a contradiction.

\subsection{TRK}
Back to the end of Case B,C,D \ref{CaseBCD Track 12}.
Suppose $\fbdb$ began with $b$, and $\fbdb$ began with $o$. We assume $\fbd$ must go over an even number of letters before crossing, which necessitates that $\fbdb$ must go over an odd number of letters before crossing.

This is TRK1.

Then TRK2 is immediate.

Here if $\fbb$ were to begin with $p^\circ$ or $p^-$, then $\fbi$ and $\fbs$ will follow with $p^-$, but then one of them will immediately trace. 

On the other hand, if $\fbb$ were to begin with $p^+$, then in order to not trace, it can only do so by going in between $\fbdb$ and $\fbi$, as shown in TRK3.

Now if it continues with $r^\circ$ or $r^-$, forcing $\fbp$ to begin with $r^-$, then $\fbp$ will immediately trace after crossing. But then this means $\fbb$ has to continue with $r^+$, immediately tracing.

This is a contradiction.

This concludes the cases analysis for Track 12.